\chapter{Making It Nova}

NovaBASIC can act like a simple home-computer BASIC, but it is sitting on a
machine with useful hardware. The programs in this book should use that
hardware when it makes the game clearer or more exciting.

\section{Graphics}

Use the graphics screen for boards, maps, roads, mountains, and charts. A
larger screen means old cramped layouts can breathe. Prefer simple readable
geometry over decorative clutter.

\section{Sprites}

Use sprites for objects that move every frame: cars, players, enemies, balls,
and cursors. If an object is larger than 16 by 16, use a metasprite helper once
that library is stable and available from BASIC.

\section{Sound}

Small sound cues matter. A short note for a correct move, a lower note for a
mistake, and a short tune on the title screen can make a type-in game feel
finished without making the listing hard to understand.

\section{Expanded Memory}

Use XRAM for tables, levels, and generated data when normal BASIC memory would
make the program awkward. Keep the first versions simple, then move bulky data
out only when the game needs it.

\section{The Rule}

If a Nova feature makes a game clearer, use it. If it only makes the listing
harder to type, save it for an advanced version.
